% Глава 5. Цифры — только со статусом ✅ в docs/research/traceability.md и паспортом в docs/thesis/figures.md.
% Результаты пилота помечены \pending, пока их нет. Предварительные проверки: docs/research/research-log.md
% (2026-10-17), docs/requirements/quality-scenarios.md, docs/architecture/change-scenarios.md.
% Инциденты RQ-4: docs/thesis/materials-log.md, docs/journal/sessions/.
\chapter{Оценка и обсуждение}\label{ch:eval}

Глава сообщает, что уже измерено и что добавит пилот. Раздел~\ref{sec:eval-status} подводит итог по
каждому исследованию. Раздел~\ref{sec:eval-prelim} представляет предварительную техническую проверку.
Разделы~\ref{sec:eval-rq1}--\ref{sec:eval-rq3} описывают результаты по RQ-1--RQ-3, которые зависят от данных
пилота, раздел~\ref{sec:eval-rq4} --- уже собранные доказательства по RQ-4. Раздел~\ref{sec:eval-discussion}
обсуждает выводы, раздел~\ref{sec:eval-limits} --- их границы.

\section{Состояние оценки}\label{sec:eval-status}

\begin{table}[ht]
\centering\small
\caption{Состояние четырёх исследований.}
\label{tab:eval-status}
\begin{tabularx}{\linewidth}{@{}P{2.8cm}LL@{}}
\toprule
Исследование & Готово сейчас & Нужно \\
\midrule
1. Согласие (RQ-1) & Разметка в Telegram и пакетах; метрики по фонемам, $\kappa$, $\alpha$ считаются
автоматически & Попытки пилота; два разметчика; бутстрэп-интервалы \\
2. Ошибки (RQ-2) & Каталог ошибок; отчёты о путаницах и строках книги; анкета при входе & Данные пилота;
частоты только по подтверждённым; отдельное хранение признаков модели \\
3. Объяснения (RQ-3) & След решения; данные повторов; предрегистрированная модель & Переключатель условия
по карточкам; отложенный тест; пилот \\
4. Обвязка (RQ-4) & Журнал хуков (с 18 октября 2026~г.); журналы сессий и запросов; ADR; инциденты &
Подсчёты по всему проекту; кодирование инцидентов \\
\bottomrule
\end{tabularx}
\end{table}

\section{Предварительная техническая проверка}\label{sec:eval-prelim}

\subsection{Отличает ли конвейер верное чтение от неверного?}

До того как тренажёром воспользовался хоть один ученик, конвейер проверили на чтении эксперта ---
обучающем чтении Махмуда Халиля аль-Хусари \cite{SRC-042}. Четыре аята отправлены как попытки для их
собственного текста, три --- как попытки для другого текста. Баллы --- в таблице~\ref{tab:husary}.

\begin{table}[ht]
\centering\small
\caption{Баллы на чтении эксперта (17 октября 2026~г., настоящая модель).}
\label{tab:husary}
\begin{tabular}{@{}llr@{}}
\toprule
Условие & Аяты & Балл \\
\midrule
Аудио совпадает с текстом & 112:1 & 100 \\
 & 112:2 & 100 \\
 & 1:2 & 99 \\
 & 105:1 & 99 \\
\midrule
Аудио другого аята & 112:1 и 112:2 & 67 \\
 & 1:2 и 112:1 & 40 \\
 & 105:1 и 1:2 & 48 \\
\bottomrule
\end{tabular}
\end{table}

Верное чтение получило 99--100, несовпадающие тексты --- 40--67, далеко от порога засчитывания 80. Это
показывает, что конвейер --- преобразование, распознавание, фонетизация, выравнивание и оценка --- работает
целиком и отличает верное чтение от неверного на записи эксперта. О начинающих это пока ничего не говорит:
читает эксперт, текст коранический, пар всего семь. Синтетический арабский голос операционной системы
распознавался плохо, так что синтетическая речь не может заменить записи людей в тестах.

\subsection{Время ответа}

На первых девяти попытках с настоящей моделью на ноутбуке автора медианное время обработки составило
1,7\,с, 90-й процентиль --- 5,0\,с, что укладывается в цель 10 секунд из QAS-03. Измеряемый интервал
охватывает только распознавание и оценку; загрузка голосового, преобразование и отрисовка ответа в него не
входят, поэтому ученик ждёт дольше. В пилоте те же процентили будут сообщены по дням на сервере
развёртывания.

\subsection{Устойчивость алгоритма оценки}

Модель может вернуть любую последовательность, поэтому оценка должна вести себя разумно на любой из них.
Шесть свойств из раздела~\ref{sec:impl-verification} проверены на случайных строках фонем: 300 примеров на
свойство при каждом прогоне набора тестов и 5\,000 на свойство в разовом прогоне --- нарушений нет.

\subsection{Модифицируемость}

Добавление второго распознавателя потребовало трёх файлов и не изменило ни домен, ни сценарии, ни базу,
ни интерфейс; все архитектурные контракты и тесты прошли (QAS-10). Это подкрепляет утверждение, что
тренажёр можно использовать для сравнения моделей на одних и тех же записях, что запланировано для
докторантуры.

\section{RQ-1: согласие с преподавателями}\label{sec:eval-rq1}

\begin{notebox}
Раздел будет дописан после пилота и двойной разметки. Ниже --- планируемый вид таблиц; ни одно значение не
приводится, пока у него нет паспорта (отчёт, дата, $n$, отпечаток кода).
\end{notebox}

\begin{table}[ht]
\centering\small
\caption{Качество обнаружения ошибок на начинающих (Т-01).}
\label{tab:T-01}
\begin{tabular}{@{}lccc@{}}
\toprule
Мера & Значение & 95\,\% ДИ & $n$ \\
\midrule
Precision & \pending{Т-01} & & \\
Recall & & & \\
F1 & & & \\
Точность диагноза & & & \\
FAR (вердикт) & & & \\
FRR (вердикт) & & & \\
\bottomrule
\end{tabular}
\end{table}

\begin{table}[ht]
\centering\small
\caption{Согласие по вердиктам (Т-02).}
\label{tab:T-02}
\begin{tabular}{@{}lcc@{}}
\toprule
Пара & $\kappa$ Коэна & 95\,\% ДИ \\
\midrule
Тренажёр --- учитель A & \pending{Т-02} & \\
Тренажёр --- учитель B & & \\
Учитель A --- учитель B & & \\
\bottomrule
\end{tabular}
\end{table}

Анализ сообщит precision и recall по типам ошибок и по буквам (Р-08), $\alpha$ Криппендорфа по фонемам
между разметчиками, долю оспоренных попыток, которые разметчики исправили, и отдельно блок
самоисправлений. H1.1--H1.3 проверяются, как предрегистрировано.

\section{RQ-2: ошибки начинающих}\label{sec:eval-rq2}

\pending{результаты RQ-2}. Анализ представит частоты типов ошибок (Т-03), самые трудные строки и
минимальные пары книги (Т-04) и матрицы путаниц по всем найденным и по подтверждённым ошибкам (Р-09).
Сравнение двух матриц --- главный результат: оно показывает, какие части видимого профиля ученика создаёт
модель. H2.1 проверяется, только если есть обе языковые группы; для H2.2 нужны выходы модели по
признакам, которые тренажёр уже получает, но пока не хранит отдельно.

\section{RQ-3: эффект объяснений}\label{sec:eval-rq3}

\pending{результаты RQ-3}. До пилота тренажёру нужна ещё одна функция: переключатель между полным
объяснением и сигналом «верно / неверно» для каждой пары «ученик --- карточка» с записью условия в попытку.
Анализ сообщит изменение балла после каждого условия (Р-11), повторяемость названных ошибок и результаты
отложенного теста (Т-05) по смешанной модели из раздела~\ref{sec:method-s3}. Без контрольного условия
улучшение при повторе смешивало бы эффект объяснения с эффектом простой повторной попытки; поэтому
описательная статистика повторов, собранная до появления переключателя, доказательством не считается.

\section{RQ-4: обвязка разработки}\label{sec:eval-rq4}

Количественные меры --- блокировки по причинам (Р-13), вмешательства за сессию и итерации до зелёных
проверок (Т-08) --- будут подсчитаны по всему проекту в его конце; журнал хуков ведётся с 18 октября
2026~г. \pending{подсчёты RQ-4}. В журналах уже есть инциденты, показывающие, что обвязка поймала, что
пропустила и где понадобился человек. Таблица~\ref{tab:incidents} группирует их по виду сбоя.

\begin{table}[ht]
\centering\small
\caption{Задокументированные инциденты из журналов сессий (октябрь 2026~г.).}
\label{tab:incidents}
\begin{tabularx}{\linewidth}{@{}P{3.0cm}LP{2.8cm}@{}}
\toprule
Вид сбоя & Инцидент & Чем пойман \\
\midrule
Подгонка тестов & В принятом тесте было неверное ожидание, написанное агентом; изменение прошло только
через фазу тестов с согласия автора. & Замороженные тесты, стоп-сигнал \\
Размывание архитектуры & Новый исследовательский модуль импортировал слой запуска; точку входа
перенесли. & import-linter \\
Утечка секретов & Регулярное выражение агента не маскировало токен бота внутри URL. & Тест логов \\
Тихий дефект & Даты в SQLite теряли часовой пояс; добавлен тип столбца с UTC. & Интеграционный тест \\
Разрушительная правка & Убирая клавиатуру, агент обрезал файл вместе с клавиатурами согласия. & Проверка
типов \\
Невысказанное допущение & Бот передан автору с заглушкой распознавания, которая засчитывает всё; автор
увидел «всё верно». & Не пойман; теперь хук напоминает агенту \\
Правило осталось прозой & Тексты кнопок согласия лежали вне защищённого файла текстов, и агент мог
изменить их без подтверждения. & Ревью; перенесены, покрыты тестом \\
Проверка со слепым пятном & Правило подтверждения закрывало правку, но не запись того же файла. & Ревью \\
Проверка со слепым пятном & Шаблон в \texttt{.gitignore} скрывал пять файлов документации от всех
коммитов. & Сверка документов \\
Проверка со слепым пятном & Аудит зависимостей не видел необязательное дополнение с моделью и его
уязвимости. & Проверка зависимостей \\
Зелёное, но сломанное & Контейнер не мог писать логи, а задаче CI не хватало тестовой зависимости; оба
прошли все локальные проверки. & Ревьюер направления (СТОП) \\
Текст не соответствует действию & Кнопка обещала алфавит, а урок начинался с \emph{басмалы}. & Консольный
«водитель бота» \\
\bottomrule
\end{tabularx}
\end{table}

Уже видны три наблюдения. Первое: детерминированные проверки поймали классические сбои агента ---
подгонка тестов заблокирована и решена через автора, нарушения архитектуры и ошибки типов остановили
агента. Второе: несколько проверок были зелёными, но слепыми: правило конфигурации, шаблон исключения и
аудит зависимостей покрывали меньше, чем казалось. Каждый случай нашла другая проверка или ревью, что
говорит в пользу проверки самих проверок. Третье: самый дорогой инцидент --- заглушка распознавания --- был
сбоем коммуникации, а не вычисления; его исправили, добавив хук, который превращает напоминание в
механизм. Количественный анализ покажет, как часто такие инциденты случаются относительно объёма работы.

\section{Обсуждение}\label{sec:eval-discussion}

\paragraph{Две вероятностные системы под одной дисциплиной} Работа имеет дело с двумя системами, которые
могут пройти проверку, не сделав правильного. Агент-разработчик может добиться прохождения тестов, изменив
их; распознаватель может отклонить верное чтение или принять неверное. Средства имеют одинаковую форму:
детерминированные внешние проверки, которые система не может править, журнал каждого решения и путь
эскалации к человеку. Для агента это замороженные тесты, контракты слоёв и вердикты СТОП; для модели ---
след решения, кнопка «Бот ошибся», разметка учителями и, как следующий шаг, исход «не уверен» на основе
уверенности модели.

\paragraph{На что тренажёр может претендовать} Литература и предварительные проверки поддерживают скромную
роль: терпеливый партнёр по практике, который между уроками даёт немедленную, явную первичную обратную
связь. Утверждение, что он заменяет учителя, не подкреплено данными и противоречит самой книге. Дизайн это
отражает: оспорить вердикт легко, оспоренные попытки первыми идут к учителям, и сохраняется всё, что нужно
учителю для проверки вердикта.

\paragraph{Ожидающие выводы} \pending{интерпретация результатов RQ-1--RQ-3}.

\section{Ограничения}\label{sec:eval-limits}

Пилот небольшой и закрытый, поэтому его результаты описывают участников, а не всех русскоязычных
начинающих. Разметчики видят ответ тренажёра, что может завышать согласие. Модель обучена на чтении
Корана, а с.~10--37 книги содержат слоги и слова не из Корана; поэтому результаты сообщаются по типу
материала. Мера сходства, используемая для вердикта, груба на длинных строках. В кейс-стади обвязки один
разработчик, один проект, и автор сам же оценивает; оно иллюстративное, а не сравнительное. Наконец, текст
программы, расшифрованный со скана, ещё ждёт полной сверки человеком; каждая попытка хранит версию
программы, поэтому попытки, затронутые будущим исправлением, можно найти и исключить.
